\section{Introduction}\label{sec:intro}

Continuous (implicit) representations decode a signal at arbitrary
coordinates, and a large body of work controls \emph{which frequencies} such a
representation can express: the frequency scale of sinusoidal activations
\citep{sitzmann2020siren,liu2024finer}, the bandwidth of Fourier features
\citep{tancik2020fourier}, target-spectrum-driven hyper-parameter selection
\citep{kania2025fresh,winner2025,sec2026}, per-band capacity allocation
\citep{wieninr}, or explicit spectral priors such as Gaussian processes and
kriging \citep{rasmussen2006gp}. We refer to all of these as \emph{spectral
adaptation}. Its empirical record is uneven. In multi-frame
super-resolution, choosing the Fourier-feature bandwidth well is decisive
\citep{superf2026}; in sparse-pixel reconstruction, a matched spectrum is
worth several decibels (\S\ref{sec:exp-b}); but in single-image
arbitrary-scale super-resolution (ASISR) the same lever barely moves.

ASISR is the setting this paper grew out of. Starting from LIIF
\citep{chen2021liif}, a line of decoders---LTE's dominant-frequency
estimator \citep{lee2022lte}, our scale-consistent SC-INR with an analytic
footprint (sinc) readout, and rotation-equivariant decoders (Rot-E)
\TODO{cite Rot-E}---each injected some form of spectral or geometric
structure into the decoder. On the shared EDSR-baseline/DIV2K protocol,
three years of such designs moved Urban100 $\times4$ by about 0.35\,dB
(MetaSR 26.05\,dB \citep{hu2019metasr} vs.\ CiaoSR 26.58\,dB \citep{cao2023ciaosr}, as reported by the original papers) \src{LIT-B} and a
controlled re-training study finds the best method only 0.035\,dB ahead of the
second \citep{nasir2026empirical}. Our own controlled arms reproduce the
picture at the level of frequency bands: at $\times4$, four decoders recover
the band one octave above the LR Nyquist limit to within 0.02 of each other
\src{OUT7.1/h6c}. The natural question is \emph{why}: is the
decoder the wrong place to adapt the spectrum, or is there nothing left to
adapt?

We argue that the answer is determined by an information budget that can be
computed before any training. For a known linear observation $y=Ax+n$, the
fraction of each band that \emph{any} stationary-kernel linear estimator can
recover is bounded by the LMMSE recoverability $\Rlin$ under a generic
$1/f^2$ prior. This number depends only on $A$, the noise level and the
generic prior, and is cheap to compute. Combined with a closed-form cost of
using a wrong spectrum, it yields a three-regime map that says, per task and
per band, whether spectral adaptation can help, by how much at most, and in
which direction an error is dangerous.

\paragraph{Contributions.}
\begin{enumerate}
\item \textbf{A training-free diagnostic.} We cast spectral adaptation as the
choice of a prior spectrum (Prop.~\ref{prop:bochner}) and use the per-band
linear recoverability $\Rlin$ (Prop.~\ref{prop:bound}) to draw an a-priori
three-regime map for single-image, sparse-pixel, multi-frame and 1D signal
tasks (\S\ref{sec:zones}, Fig.~\ref{fig:map}).
\item \textbf{A closed-form mismatch cost.} Prop.~\ref{prop:mismatch} gives
the excess error of a prior that is off by a factor $\kappa$ in a band of
recoverability $R$. Its asymmetry yields a per-band tendency---\emph{wide} is safe as $R\to1$,
\emph{matching} matters at intermediate $R$, \emph{conservative} is safe where
$R$ is low---and its energy-weighted sum over bands ranks representations
(though it does not reliably pick the single best); it also explains why grid fields collapse under noise and
why large-$\sigma$ Fourier features fail (\S\ref{sec:zones}).
\item \textbf{Pre-registered cross-task validation} (\S\ref{sec:exp}).
The mismatch-cost shape and asymmetry hold on 60 task--band--$\kappa$ cells
(Spearman 0.90) \src{E2}; the zones move with the ASISR scale factor as
predicted (Spearman 0.952 over 27 cells) \src{E3}; in the low-$R$
regime trained ASISR networks exceed the linear oracle in every cell while
differing among themselves by at most 0.03 in 10 of 12 cells \src{E3},
which accounts for decoder saturation; and the benefit of the cell/footprint
readout is zero in the high-$R$ band and appears only where the network has
partially recovered the band---consistently in LTE and SC-INR
(cross-architecture Spearman 0.985, seed 1) \src{F1}; averaging the full model
over training scales predicts the ranking of these losses after the fact
(Spearman 0.82; correlational) \src{L1L2}.
\end{enumerate}

\paragraph{Scope and honesty.} Props.~\ref{prop:bound}--\ref{prop:mismatch}
follow directly from classical Wiener/LMMSE theory; our claim is not new
estimation theory but its use as an a-priori diagnostic and prescription for
spectral adaptation in continuous representations, together with a
pre-registered test across tasks. Several pre-registered criteria failed; we
report them in the main text and in \S\ref{sec:limits}.
